\documentclass[12pt]{article}
\usepackage{amsmath,amssymb,amsfonts}
\usepackage{tikz}
\usepackage[margin=0.5in]{geometry}
\newcommand{\pdata}{p_\mathrm{data}}
\begin{document}
\begin{center}
Name: \rule{3in}{0.4pt} \hfill Total: 16 points
\end{center}
We will consider a 2D Ornstein-Uhlenbeck process given by the following SDE:
		\begin{equation}
			dX_t = -\theta\: X_t \, dt + \sigma\: dW_t,
		\end{equation}
		where $\theta, \sigma > 0$ are constants. We will treat the paths $\{X_t\}$ as the solution of a random dynamical system, where the randomness comes from the standard Wiener process, $W_t,$ in 2D. There is a unique, invariant measure of the process given by
		\begin{equation}
			X_\infty \sim \mathcal{N}\!\left(0, \frac{\sigma^2}{2\theta} I_2\right).
		\end{equation}
We can use Ito integrals to obtain that the process $X_t$ is the following random variable,
\begin{equation}
	X_t = e^{-\theta \: t} X_0 + \sigma \int_0^t e^{-\theta(t-s)} dW_s.
\end{equation}

\begin{enumerate}

	\item Below we see a simulation of the densities (estimated) of sample paths for a fixed Wiener path. That is, for a fixed path $W$, we see the distribution $\varphi^t_{W\sharp} \rho_0$ for some initial density, $\rho_0$ (a shifted Gaussian), at different times $t.$  As $t$ increases, the densities will concentrate at a single point. True or False? Justify your answer. (2 points)   
		\begin{figure}[h]
			\includegraphics[width=\textwidth]{ou_simulation_499.png}
		\end{figure}
	
\newpage
\item What is the distribution of the limiting points/set of points obtained from such simulations with fixed Wiener paths? In other words, what is $\lim_{t\to\infty}\mathbb{E}_{W \sim \mathbb{P}} \varphi^t_{W\sharp} \rho_0,$ where $\mathbb{P}$ is the Wiener measure? (2 points)

\vfill

\item Let $P_t(x, x_0)$ be the transition probability density of $X_t = x$ conditioned on $X_0 = x_0$. The Koopman operator $U^t:L^2(\mu) \to L^2(\mu)$ is given by 
	\begin{equation}
		U^t\: f(x_0) = \int P_t(x,x_0)\: f(x)\: dx. 
	\end{equation}	
	Relate this definition to the Koopman operator for the deterministic system, $\varphi^t_W$. (2 points)

\vfill
\newpage

\item What is the Frobenius-Perron operator associated with the Koopman semigroup $U^t$? Write down how it acts on probability densities. \emph{Hint: it is the $L^2(dx)$ adjoint of $U^t$.} (2 points)


\vfill

\item Let $H_n(x) = h_{n_1}(x^{(1)})\: h_{n_2}(x^{(2)})$ be a product of 1D Hermite polynomials. Here, $n = (n_1, n_2)$ is a tuple of non-negative integers and $(x^{(1)}, x^{(2)})$ are coordinates of a point $x.$
These are eigenfunctions of $U^t$ with eigenvalues $e^{-\theta (n_1 + n_2) t}.$ Knowing this, what do you expect the Koopman approximation to look like from EDMD, when choosing the span of $H_n$ as the basis? (3 points)
\vfill


\item Briefly describe how you would perform EDMD for this system using the basis functions from the previous part. (3 points)
\vfill
\newpage
\item Any initial density converges to a Gaussian exponentially fast. Argue why or why not. (2 points)
\end{enumerate}



\end{document}
